\section{Classicality and descent}\label{sec:classicality}

Theorem~\ref{thm:potential-promodularity} places the target representation
in the closed Hecke support after a solvable totally real base change.
We now produce a locally algebraic vector with its eigensystem.
The additional work is needed when the local representation at~$2$ is
reducible: the two possible orders in a unitary principal series need
not have the same locally algebraic vectors.
Throughout this section the target $r$ is not virtually solvable, as
arranged in the preliminary reductions.  In particular, its restriction
to every finite extension is absolutely irreducible.

We retain the conventions of Section~\ref{sec:completed}:
$\chi=\det r$, the central character of completed forms is
$\psi=\chi\eps$, and the local block parameter is $r\otimes\eps$.
Cyclotomic exponents are counted with $\operatorname{exp}(\eps)=1$.
Thus a de Rham character of exponent $a$ has Hodge--Tate weight $-a$
in the convention used by the ordinary modularity theorem below.
The field $F$ will always be totally real, solvable over~$\QQ$, of even
degree, and split completely at~$2$.
Until the final subsection, assume that $r|_{\GQtwo}$ is reducible.
We will retain its higher-exponent line in a global family, show that
the joint trace-and-character support is finite over weight space, and
use dense points that become modular after normalization and dualization.
Those points determine the inducing order to retain at the target.
Specialization will then give an actual vector, whose local algebraicity
is proved directly.  The locally absolutely irreducible case will use
Proposition~\ref{prop:block-comparison} in the final subsection.

\subsection{The local line and its global family}

We first record the rank-one consequence of Bloch--Kato that fixes the
orientation.  For a de Rham character $\theta$ of $G_{\QQ_2}$, write
$H^1_g(\QQ_2,\theta)$ for the subspace of de Rham extension classes.
The local Euler characteristic, duality, and the Bloch--Kato formulas give
\begin{align}
 \dim H^1(\QQ_2,\theta)
 &=1+\dim H^0(\QQ_2,\theta)
       +\dim H^0(\QQ_2,\theta^{-1}\eps),\label{eq:rank-one-euler}\\
 \dim H^1_g(\QQ_2,\theta)
 &=\dim D_{\mathrm{dR}}(\theta)/\Fil^0
   +\dim H^0(\QQ_2,\theta)
   +\dim D_{\mathrm{cris}}(\theta^{-1}\eps)^{\varphi=1}.
   \label{eq:rank-one-geometric}
\end{align}
The second formula follows from the exact annihilators of
\cite[Proposition~3.8]{BlochKato1990}, the $H^1_f$ and $H^1_e$
dimension formulas of \cite[Corollary~3.8.4]{BlochKato1990}, and local
duality; it is the rank-one consequence used here.
These formulas have no restriction excluding~$2$.

\begin{lemma}\label{lem:higher-line}
If $r|_{\GQtwo}$ is reducible and has distinct cyclotomic exponents,
it has an invariant line whose exponent is greater than that of the
quotient.
\end{lemma}
\begin{proof}
Characters occurring as subquotients of a de Rham representation are
de Rham.  Suppose a chosen subcharacter has smaller exponent than the
quotient, and let $\theta$ be their ratio.  Its exponent is strictly
negative, so $D_{\mathrm{dR}}(\theta)/\Fil^0=0$ and
$H^0(\QQ_2,\theta)=0$.  The character $\theta^{-1}\eps$ has strictly
positive exponent.  If its finite inertial part is nontrivial, its
crystalline module is zero; otherwise its crystalline Frobenius
eigenvalue has nonzero valuation and cannot be~$1$.
Equation~\eqref{eq:rank-one-geometric} therefore gives $H^1_g=0$.
The extension splits, and its other summand is the required line.
\end{proof}

\begin{lemma}\label{lem:ordinary-character-family}
Suppose $r|_{\GQtwo}$ is reducible, and record the line of
Lemma~\ref{lem:higher-line}.  Fix $\chi=\det r$.
After enlarging coefficients, there is a
complete local integral $\OO$-algebra $\mathscr D$, together with its
fixed global determinant $\chi$, its global trace, and a recorded local subcharacter
$\eta_1$, with the following properties.
\begin{enumerate}
\item $\mathscr D$ is horizontal, contains the target point, and
      $\dim\mathscr D\ge2$.  Put $\eta_2=\chi\eta_1^{-1}$.
      On its globally absolutely irreducible locus the local
      representation has subcharacter $\eta_1$ and quotient $\eta_2$.
\item For one field $F$ as above, the restricted trace lies in the
      Hecke support of completed forms at a fixed tame level.
\item $\mathscr D$ is finite over the closed image of the global
      pseudodeformation ring for $F$.  It is also finite over the
      image of the product of the dyadic pseudodeformation rings.
\end{enumerate}
\end{lemma}
\begin{proof}
Use a framed global deformation chart with fixed determinant and an
invariant line at~$2$, in an integral basis for the target and its line.
Consider its completed local ring at the characteristic-zero target
point.  Corollary~\ref{cor:charzero-line-dimension} gives dimension
at least five for the local line chart and dimension two for the odd
real factor, over the characteristic-zero coefficient field.
The relative presentation of Proposition~\ref{prop:relative-presentation}
has generators minus relations at least $-3$: its dual invariant term
vanishes because the global representation is not virtually solvable.
The global completed local ring therefore has dimension at least
$5+2-3=4$.  The characteristic-zero dimension comparison in
Lemma~\ref{lem:curve-completion} then gives a horizontal integral
component through the target in the original chart of dimension at
least five.

Contract that component to its global trace and local subcharacter
values.  On the absolutely irreducible locus the only positive
dimensional fibers are conjugacy fibers, of dimension three.  The local
characters are distinct on a dense open set, so recording the character
does not leave an extra moving-line fiber there.  The dimension
comparison of Lemma~\ref{lem:dimension-comparison} gives a domain
$\mathscr D$ of dimension at least two.  Properness of the invariant-line
incidence preserves the recorded subcharacter condition on the
irreducible locus.  Theorem~\ref{thm:potential-promodularity}, applied to
an unrestricted framed component containing this family, gives one
common $F$ and one tame level with the asserted trace support.

We give the finiteness argument, including the recorded character.
Write $T_{\mathscr D}$ for the trace of the family.
Let $A\subset\mathscr D$ be the closed restricted trace image, and let
$q=[G_{\QQ}:G_F]$.  For every $g\in G_{\QQ}$, $g^q\in G_F$ after
replacing $q$ by the order of the finite quotient if necessary.
The identity
\[
 T_{\mathscr D}(g^q)=P_q\bigl(T_{\mathscr D}(g),\chi(g)\bigr)
\]
has $P_q(X,Y)$ monic of degree $q$ in~$X$.  Hence every downstairs
trace is integral over~$A$.  A local character value satisfies
\[
 \eta_1(g)^2-T_{\mathscr D}(g)\eta_1(g)+\chi(g)=0,
 \qquad g\in G_{\QQ_2},
\]
and is integral as well.  The global determinant algebra has finitely
many topological trace generators, and a continuous local character is
determined topologically by its values on the uniformizer, $-1$, and~$5$.
Adjoining these finitely many integral elements to~$A$ gives a finite
complete algebra, hence a closed subalgebra of~$\mathscr D$; it contains
its topological generators and is all of~$\mathscr D$.
Finally, Proposition~\ref{prop:block-comparison} makes the Hecke algebra
finite over the product of the local central pseudodeformation images.
Transitivity of finiteness proves the last assertion.
\end{proof}

We may discard the proper globally reducible loci, including reducibility
after restriction to~$F$.  We may also require
\begin{equation}\label{eq:generic-ordinary-ratios}
 \eta_1\eta_2^{-1}\notin\{1,\eps,\eps^{-1}\}.
\end{equation}
Indeed, each exceptional equality and the fixed determinant imply
$\eta_1^2=\chi\eps^a$ for $a\in\{0,1,-1\}$.  The local trace image
then has dimension at most one: its character values are integral over
the fixed coefficient ring.  Because $2$ splits in~$F$ and the family
comes from~$\QQ$, the same holds at all dyadic factors.  Finiteness in
Lemma~\ref{lem:ordinary-character-family} shows that none of these loci
contains $\Spec\mathscr D$.  Characteristic-zero points with finite
coefficient fields satisfying these conditions are consequently dense.

\subsection{The ordinary torsion module}

We now have a family carrying both global traces and the chosen local
character.  To control its character parameter, we construct a torsion
space on which the traces and torus act jointly.  Its finiteness over
torus weight space will constrain the dimension of $\mathscr D$.

Fix the Hecke summand of the completed forms over~$F$ carrying the
restricted family, and write
$\mathcal H$ for its torsion function space with coefficients
$\EE/\OO$.  Put
\[
 G_2=\prod_{v\mid2}\GL_2(\QQ_2),\quad
 N_0^-=\prod_{v\mid2}n^-(\ZZ_2),\quad
 T_1=\prod_{v\mid2}(1+4\ZZ_2)^2,
\]
where
\[
 n^-(a)=\begin{pmatrix}1&0\\a&1\end{pmatrix},\qquad
 n^+(x)=\begin{pmatrix}1&x\\0&1\end{pmatrix},\qquad
 \mathbf p=\bigl(\operatorname{diag}(1,2)\bigr)_{v\mid2}.
\]
All matrix formulas involving several places are interpreted
coordinatewise.  On $\mathcal H^{N_0^-}$ define
\begin{equation}\label{eq:ordinary-U}
 U=\sum_{a\in(\ZZ/2\ZZ)^{[F:\QQ]}}n^-(a)\mathbf p.
\end{equation}
The actions are the right translation actions with the representation
group law.  Let
\[
 W=\{f\in\mathcal H^{N_0^-}:U^Nf=f\text{ for some }N>0\},
 \qquad M=W^\vee,
 \qquad \Omega=\OO[[T_1]].
\]
The dual is the compact Pontryagin dual.  We transpose the commuting
operator actions without inversion, so their scalar eigenvalues retain
the conventions above.

\begin{proposition}[Ordinary weight finiteness]
\label{prop:ordinary-weight-finiteness}
The space $W$ is an $\OO$-submodule stable under $T_1$, the tame Hecke
operators, and $U^{\pm1}$, and these actions commute.  Its dual $M$ is
finite over~$\Omega$.  The closed commutative algebra generated by these
actions on~$M$ is finite over~$\Omega$.
\end{proposition}
\begin{proof}
In one factor the identity
\begin{equation}\label{eq:U-representatives}
 n^-(a+2b)\operatorname{diag}(1,2)
 =n^-(a)\operatorname{diag}(1,2)n^-(b)
\end{equation}
proves independence of representatives and preservation of lower
invariance.  Conjugation by $t=\operatorname{diag}(u_1,u_2)\in T_1$
replaces $a$ by $(u_2/u_1)a$, which permutes the residue classes.
Thus $tU=Ut$ on the lower-invariant space.  The tame Hecke actions commute
with these local actions.  A common multiple of two periods is a period
of the sum of the corresponding vectors.  Finally $U$ is bijective
on~$W$: on a vector of period $N$, its inverse is $U^{N-1}$, independently
of the chosen period.

We show that one compact open subgroup fixes $W[\varpi]^{T_1}$.
For $j\ge2$, $x\in2^j\ZZ_2$, and $a\in\ZZ_2$, direct multiplication
gives
\begin{align}
 n^+(x)n^-(a)\operatorname{diag}(1,2)
 ={}&n^-\!\left(\frac{a}{1+xa}\right)
       \operatorname{diag}(1,2)\notag\\
 &\cdot\operatorname{diag}(1+xa,(1+xa)^{-1})
       n^+\!\left(\frac{2x}{1+xa}\right).
       \label{eq:U-upper-improvement}
\end{align}
If $f$ is $T_1$-fixed and upper-$2^{j+1}$-invariant, the final two
factors fix~$f$.  The map $a\mapsto a/(1+xa)$ permutes the lower
residue classes, and \eqref{eq:U-representatives} absorbs the changes
of lift.  Thus $Uf$ is upper-$2^j$-invariant.  Once upper-$4$ invariance
is reached, it is preserved by the same identity.

Each smooth vector of $W[\varpi]^{T_1}$ starts at some finite upper
level.  Iterate the preceding improvement far enough, taking an
iteration number divisible by its period.  The resulting vector is the
original vector, so it is fixed by the single open compact subgroup
generated by $N_0^-$, $T_1$, and $\prod_{v\mid2}n^+(4\ZZ_2)$.
Proposition~\ref{prop:completed-forms} gives smooth admissibility;
therefore $W[\varpi]^{T_1}$ is finite-dimensional over~$\kk$.
Compact Nakayama now gives finite generation of~$M$ over~$\Omega$.

Since $\Omega$ is noetherian and $M$ is finite, $\End_\Omega(M)$ is
finite over~$\Omega$.  The closed algebra in the statement is an
$\Omega$-submodule of this finite module, and hence is finite too.
\end{proof}

Let $R$ denote the fixed-determinant global pseudodeformation ring
over~$F$.  The trace and torus actions define
\[
 Q=R\widehat\otimes_{\OO}\Omega
       \longrightarrow\End_\Omega(M).
\]
Write $S$ for its image.  It is finite over~$\Omega$ and is closed:
the source is compact and the target is Hausdorff.  The larger action
algebra in Proposition~\ref{prop:ordinary-weight-finiteness}, denoted
by~$\mathcal A$, is finite over~$S$ and contains both $U$ and $U^{-1}$.

\subsection{Finite joint images and the weight map}

Set
\[
 \Lambda=\OO[[1+4\ZZ_2]]\simeq\OO[[T]],\qquad
 [5]=1+T.
\]
The map $\Lambda\to\mathscr D$ sends $[5]$ to $\eta_1(5)$, using local
reciprocity.  At each place of~$F$ above~$2$, there are two possible
inducing assignments
\begin{equation}\label{eq:two-ordinary-orders}
 (\eta_1\eps,\eta_2),\qquad (\eta_2\eps,\eta_1).
\end{equation}
For each choice $s$ of an assignment at every dyadic place, the fixed
determinant and these formulas define a homomorphism
$s:\Omega\to\Lambda$.  Only characters on $1+4\ZZ_2$ enter this map;
the unramified character values remain in the trace algebra.

\begin{lemma}[Finite joint-image argument]\label{lem:joint-image}
Let $\mathscr D$ be the domain of
Lemma~\ref{lem:ordinary-character-family}, let $M=W^\vee$ be the
ordinary module above, and use the weight map
$\Lambda\to\mathscr D$ defined by $[5]\mapsto\eta_1(5)$.
Suppose a dense set of finite-coefficient characteristic-zero points
of~$\mathscr D$ occurs in the joint trace-and-weight support of~$M$,
with one of the finitely many assignments~$s$ in
\eqref{eq:two-ordinary-orders} at each point.  Then
$\mathscr D$ is finite and dominant over~$\Lambda$; in particular,
$\dim\mathscr D=2$.
\end{lemma}
\begin{proof}
For each $s$, form $S_s=S\widehat\otimes_{\Omega,s}\Lambda$, and let
$T_s$ be the image of $R\widehat\otimes_{\OO}\Lambda$ in~$S_s$.
Both are finite $\Lambda$-algebras, and their images are closed.  Let
$I_s$ be the kernel of the latter map.  Each point in the dense set
kills some $I_s$ after mapping to~$\mathscr D$.  Hence
\[
 \Spec\mathscr D=\bigcup_s V(I_s\mathscr D).
\]
The spectrum is irreducible, so one of these finitely many closed sets
is the whole spectrum.  Equivalently, the prime kernel of
$R\widehat\otimes\Lambda\to\mathscr D$ contains one~$I_s$.
The closed trace-plus-weight image $B$ in~$\mathscr D$ is therefore a
quotient of~$T_s$ and is finite over~$\Lambda$.

By Lemma~\ref{lem:ordinary-character-family}, $\mathscr D$ is finite
over the restricted trace image $A$, and $A\subset B$.  The same finite
set of $A$-module generators generates $\mathscr D$ as a $B$-module.
Thus $\mathscr D$ is finite over~$\Lambda$.  Its dimension is at least
two, whereas $\dim\Lambda=2$; a nonzero prime kernel in the domain
$\Lambda$ would lower that dimension.  The map is therefore injective,
and the dimension is exactly two.
\end{proof}

The hypotheses of the lemma hold here.  At each finite-coefficient
characteristic-zero point outside the excluded loci,
Proposition~\ref{prop:block-comparison}
supplies a Banach subrepresentation of completed forms whose local
factors are unitary continuous inductions in one of the orders
\eqref{eq:two-ordinary-orders}; the two possibilities are precisely
those of \cite[Corollary~6.10(3)]{PaskunasTung2021}
for the normalized parameter $r\otimes\eps$.
In $\Ind_B^{\GL_2(\QQ_2)}(\delta_1\otimes\delta_2)$, with $B$ the
upper Borel and unnormalized induction, take the function supported
on $Bn^-(\ZZ_2)$ and equal to~$1$ on $n^-(\ZZ_2)$.
It is lower-invariant and has torus character $(\delta_1,\delta_2)$.
Subdivision into the two lower disks shows that its $U$-eigenvalue is
$\delta_2(2)$, a unit.  The product vector has the analogous properties
at all dyadic places.

Its nonzero image in completed forms can be scaled into the integral
lattice.  Modulo every power of the coefficient uniformizer its unit
eigenvalue has finite multiplicative order, so the reduction belongs
to~$W$.  Duality consequently places this trace-and-weight point in
$\Spec S$.  Finite coefficient extension is harmless, either by flat
extension of the torsion space or by taking scalar coordinates.
Lemma~\ref{lem:joint-image} applies.

Thus the family is finite and dominant over the one-variable weight
space.  We can now choose a dense set of character values there and
lift them to the family.  At these points the local extension will
satisfy the ordinary modularity theorem after the normalization and
dualization described next.

\subsection{A dense set of ordinary modular points}

Let $w=\operatorname{exp}(\chi)$.  After a finite extension of~$\EE$,
choose $c$ with $c^2=5^{w+1}$.  Consider weights
\begin{equation}\label{eq:translated-torsion-weights}
                 \eta_1(5)=c\zeta,
       \qquad \zeta\in\mu_{2^\infty}.
\end{equation}
These are continuous characters of $1+4\ZZ_2$: this group is procyclic
with generator~$5$, and $c\zeta$ lies in the open unit neighborhood
of~$1$ in its finite coefficient field.  We are specifying a character
value, not asserting that $(w+1)/2$ is an integral Hodge--Tate weight.

The weights in \eqref{eq:translated-torsion-weights} are Zariski dense
in~$\Spec\Lambda$.  Indeed, after dividing by its maximal scalar power,
a nonzero element of $\OO[[T]]$ is a unit times a distinguished
polynomial by Weierstrass preparation and has only finitely many zeros
in the open unit disk.  Finite dominance then gives a dense set of
coefficient points of~$\mathscr D$ above these weights, even after
removing any prescribed proper closed subset.  To see the last assertion
directly, a nonzero ideal of the finite integral domain $\mathscr D$
meets~$\Lambda$ nontrivially: an integral equation of minimal degree for
a nonzero element of the ideal has nonzero constant term.  Thus every
proper closed subset has proper closed image in~$\Spec\Lambda$.

On inertia, the ratio at these points has the form
\begin{equation}\label{eq:ordinary-ratio-xi}
                    \eta_1\eta_2^{-1}=\eps\xi,
\end{equation}
where $\xi$ has finite order.  This follows from
$\eta_1^2/\chi$ on $1+4\ZZ_2$, and the remaining unit quotient is
finite.  The finite part of $\chi$ is fixed; at most finitely many
choices of~$\zeta$ make $\xi$ trivial.  We discard them, as well as the
globally reducible loci before and after restriction to~$F$.

There is a single continuous normalizing character over the family.
The character $\eta_2$ on inertia factors through the local unit group
$\ZZ_2^\times$ by local class field theory.  Denote this unit character
by $\eta_{2,\mathrm{u}}$, and put
\begin{equation}\label{eq:normalizing-character}
 \lambda(g)=\eta_{2,\mathrm{u}}\bigl(\eps(g)\bigr)^{-1},
 \qquad g\in G_{\QQ}.
\end{equation}
This character is continuous and unramified outside~$2$.  At a point
chosen above, write $r'=r_x\otimes\lambda_x$ and
$\eta'_j=\eta_{j,x}\lambda_x$.  Then
\begin{equation}\label{eq:normalized-ordinary-characters}
                \eta'_1=u_h\xi\eps,
          \qquad\eta'_2=u_l,
\end{equation}
with $u_h,u_l$ unramified and $\xi$ of nontrivial finite inertia.

\begin{lemma}\label{lem:dense-Thorne-points}
Every representation $r'$ just constructed is odd, absolutely
irreducible, finitely ramified, ordinary, and potentially crystalline
with cyclotomic exponents $\{1,0\}$.  It is modular.
\end{lemma}
\begin{proof}
The two diagonal characters in
\eqref{eq:normalized-ordinary-characters} are de Rham.  Their ratio
$\theta=(u_h/u_l)\xi\eps$ has
$H^0(\theta)=H^0(\theta^{-1}\eps)=0$, the latter because the finite
inertial character is nontrivial.  Thus \eqref{eq:rank-one-euler} gives
$\dim H^1(\QQ_2,\theta)=1$.  In
\eqref{eq:rank-one-geometric} the de Rham tangent term has dimension
one, while the other two terms vanish.  Hence every extension class
belongs to $H^1_g$ and $r'$ is de Rham.

The $p$-adic monodromy theorem makes it potentially semistable
\cite[Theorems~0.6 and~5.19]{Berger2002}.  Its two finite inertial
characters are distinct.  Monodromy commutes with inertia, so it
preserves their two eigenspaces; it vanishes on each by nilpotence.
Thus it is potentially crystalline.

Twisting preserves absolute irreducibility and finite ramification.
For a complex conjugation $c_\infty$, $\lambda(c_\infty)^2=1$, so
the determinant remains odd.  Finally, in the convention
$\HT(\eps)=-1$, the dual has weights $\{0,1\}$ and ordinary shape
\[
 (r')^\vee|_{\GQtwo}\sim
 \begin{pmatrix}
   u_l^{-1}&*\\0&\eps^{-1}u_h^{-1}\xi^{-1}
 \end{pmatrix}.
\]
All the hypotheses of \cite[Theorem~D]{Thorne2026} are now satisfied.
That theorem makes the dual modular, and hence makes $r'$ modular.
\end{proof}

The modular representations in this lemma have a common bound for
their tame conductors.  Indeed, the original noetherian Galois family
has, at each non-dyadic bad prime, a common open inertia subgroup on
which its action is unipotent, as in
Lemma~\ref{lem:inertia-basechange}.  The finite inertial quotient is
fixed, and the rank-two monodromy contribution is bounded.  This bounds
the Artin conductor after the fixed base change to~$F$.
The normalizing characters add no non-dyadic inertia.  Consequently,
after replacing the tame compact subgroup by one smaller subgroup,
all these modular forms occur at a single tame level.  Solvable base
change, iterated through prime-degree cyclic extensions, and
Jacquet--Langlands put them in the definite forms over~$F$;
see \cite[Chapter~11, Lemma~11.3 and Proposition~11.4]{Langlands1980},
\cite[Chapter~III, Theorems~4.2 and~5.1]{ArthurClozel1989}, and
\cite[Theorem~16.1]{JacquetLanglands1970}.

\subsection{The required inducing order}

We next track the order at the dense modular points, rather than
selecting it from a semisimple Banach fiber.  In the geometric
reciprocity convention $\eps=x|x|$, one has $\WD(\eps)=|\cdot|$.
For the normalized characters \eqref{eq:normalized-ordinary-characters},
the two characters of $\WD(r'\otimes\eps)$ are
\[
              \theta_h=u_h\xi|\cdot|^2,
          \qquad\theta_l=u_l|\cdot|.
\]
Here the same notation for a finite or unramified character denotes
its character under local reciprocity.  A regular algebraic Hilbert
form on $\GL_2$ is essentially self-dual via its central character.
Local--global compatibility
identifies the smooth principal series: the underlying semisimple Weil
representation has these distinct inertial characters, and no
monodromy can join them.  This also follows from the compatibility up
to the monodromy order of
\cite[Corollary~1.2]{BLGGTLocalGlobal2014}, since that order preserves
the underlying semisimple Weil representation and monodromy is zero
in this case.

In the conventions of Proposition~\ref{prop:block-comparison}, put the
smooth inducing characters in the order
\[
 (\theta_l,\theta_h|\cdot|^{-1})
             =(u_l|\cdot|,u_h\xi|\cdot|).
\]
The algebraic factor at cyclotomic exponents $\{1,0\}$ is $\det$.
Multiplication by that factor changes the pair to
\[
                 (u_l\eps,u_h\xi\eps)
                         =(\eta'_2\eps,\eta'_1).
\]
For unnormalized smooth induction, the intertwiner changes an inducing
pair $(A,B)$ to $(B|\cdot|,A|\cdot|^{-1})$.  The two orders here are
therefore intertwined with the appropriate modulus factors, with no
exceptional ratio because their inertial characters differ.
Equivalently, this is the
weight-two instance of the ordinary universal completion formula
\cite[Proposition~2.2.1]{BreuilEmerton2010}; that proposition applies
at~$2$ and concerns the locally algebraic induction.
In particular, the lower-disk vector used above already belongs to
this locally algebraic representation.

Undo the twist by multiplying definite functions by
$\lambda_x^{-1}$ evaluated on the reduced norm, via global reciprocity.
Rational reduced norms are totally positive, so this is well-defined
on the definite quotient.  It changes the central character back to
$\chi\eps$ and changes the inducing assignment to
\begin{equation}\label{eq:favorable-ordinary-order}
                       (\eta_2\eps,\eta_1)
                 \quad\text{at every }v\mid2.
\end{equation}
Lower invariance is preserved, and the $U$-eigenvalue is multiplied by
a unit.  All dyadic places have the same order because they are split
over~$\QQ$ and the family and normalizing character descend from~$\QQ$.

At this common tame level, let $M$ and $S$ again denote the ordinary
dual and its joint trace-and-weight image algebra.  Write
$s_{\mathrm{fav}}:\Omega\to\Lambda$ for the single assignment
\eqref{eq:favorable-ordinary-order}, and set
\[
 S_{\mathrm{fav}}
   =S\widehat\otimes_{\Omega,s_{\mathrm{fav}}}\Lambda,
 \qquad
 I_{\mathrm{fav}}
   =\ker\bigl(R\widehat\otimes_{\OO}\Lambda
                         \longrightarrow S_{\mathrm{fav}}\bigr).
\]
Let $T_{\mathrm{fav}}$ be the image of this map.  The untwisted vectors
just constructed show that $I_{\mathrm{fav}}$ vanishes at the dense
family points used in Lemma~\ref{lem:dense-Thorne-points}.  Its image
in the domain $\mathscr D$ is therefore zero.  In particular, the
original target defines a coefficient point of $T_{\mathrm{fav}}$,
including when its local character ratio satisfies one of the
exceptional equalities excluded from the dense testing set.

The algebra $S_{\mathrm{fav}}$ is finite over $T_{\mathrm{fav}}$,
so lying over lifts this point to $S_{\mathrm{fav}}$ after a finite
coefficient extension.  Its restriction to $S$ has the target trace
and the torus character prescribed by $s_{\mathrm{fav}}$.
Since $M$ is finite over $S$ and the $S$-action is faithful,
$\Supp_S M=\Spec S$.  Thus this point belongs to the joint
trace-and-weight support needed below.

The favorable character order has now reached the target.  To use it
in a classicality argument, we still need a nonzero vector carrying
that order.  The next step obtains one from a scalar quotient of the
finite action module, without assuming that ordinary vectors or locally
algebraic vectors commute with specialization.

\subsection{Specialization to a genuine vector}

\begin{lemma}[Eigenvector specialization]
\label{lem:eigenvector-specialization}
Let $S$ be the joint trace-and-weight image algebra acting on $M=W^\vee$
at the fixed tame level above, and let $\mathcal A$ be its finite faithful
action algebra containing $U^{\pm1}$, as in
Proposition~\ref{prop:ordinary-weight-finiteness}.  Let a
finite-coefficient characteristic-zero point of $\Spec S$ belong to
the support of $M$.  After a finite coefficient extension it gives a
nonzero vector in completed forms with the
specified trace and torus eigencharacters, invariant under~$N_0^-$,
and with a unit $U$-eigenvalue.
\end{lemma}
\begin{proof}
The algebra $\mathcal A$ containing $U^{\pm1}$ is finite over~$S$.
Since it is the faithful action algebra of~$M$, and $M$ is finite over
it, its spectrum is the support of~$M$.  Lying over the given point
and passing to a finite coefficient field $E'$ therefore gives a
joint scalar quotient of $M[1/2]$.  Write $\OO'$ for its ring of
integers, with uniformizer~$\varpi'$.

The corresponding nonzero map
\[
                    M\otimes_{\OO}\OO'\longrightarrow E'
\]
is continuous.  Indeed, $M$ is finite over weight space and its weight
values lie in the formal open unit disk.  Its image is a compact
$\OO'$-submodule spanning $E'$, hence a lattice; rescale it to~$\OO'$.
The scalar actions preserve this lattice.  In particular the
$U$-eigenvalue is a unit because $U^{-1}$ preserves it as well.

Applying $\OO'$-linear Pontryagin duality, with values in $E'/\OO'$,
produces an equivariant injection
\[
               E'/\OO'\lhook\joinrel\longrightarrow
                    W_{\OO'}\subset\mathcal H_{\OO'}.
\]
The image $f_n$ of $\varpi'^{-n}$ has exact order $\varpi'^n$ and
satisfies $\varpi'f_{n+1}=f_n$.  Taking the Tate module gives a nonzero
continuous integral completed-forms vector.  Explicitly, a compatible
sequence of continuous functions modulo $\varpi'^n$ determines a
continuous $\OO'$-valued function, because each of its reductions is
continuous; conversely every such integral function gives that
sequence.  Exact orders ensure that the resulting vector is nonzero.
All Hecke and torus eigenrelations, lower invariance, and the
$U$-eigenrelation pass to this inverse limit.  This construction finds
a divisible copy inside~$W$; it does not require $W$ itself to be
divisible.
\end{proof}

Apply the lemma to the target.  Its torus eigencharacter, restricted
to a sufficiently small open subgroup, is algebraic with exponents
$e_{1,v},e_{2,v}$ satisfying
\begin{equation}\label{eq:nonnegative-algebraic-degree}
 e_{2,v}-e_{1,v}
   =\operatorname{exp}(\eta_1)-\operatorname{exp}(\eta_2)-1\ge0.
\end{equation}
The inequality follows from the choice of the higher line and the
integrality and distinctness of the target weights.
We have therefore obtained a genuine ordinary eigenvector whose torus
exponents have the order needed for an algebraic representation.  It
remains to identify the algebraic form of its orbit on a sufficiently
small open subgroup.

\subsection{A unit ordinary vector is locally algebraic}

\begin{lemma}\label{lem:polynomial-orbit}
Let $V$ be a continuous unitary Banach representation of~$G_2$ over a
finite extension $E'/\QQ_2$.  Suppose $0\ne f\in V$ is fixed by
$N_0^-$, has torus character with integral exponents
$e_{1,v},e_{2,v}$ on a sufficiently small open subgroup, and satisfies
$Uf=\alpha f$ with $\alpha\in\OO'^\times$.  If
$\delta_v=e_{2,v}-e_{1,v}\ge0$ for all~$v$, then $f$ is locally
algebraic, of algebraic type
\[
       A=\bigotimes_{v\mid2}
                \det^{e_{1,v}}\operatorname{Sym}^{\delta_v}(E'^2).
\]
\end{lemma}
\begin{proof}
Choose $m\ge2$ so the stated torus character holds on
$\prod_v(1+2^m\ZZ_2)^2$, and put
$X=(2^m\ZZ_2)^{[F:\QQ]}$.  Lower invariance gives
\[
 U^jf=\sum_{b\bmod 2^j}n^-(b)\mathbf p^jf,
\]
independently of the integral lifts of~$b$.  For $x\in X$ the map
$a\mapsto a/(1-xa)$ permutes these residue classes.  With
$c=1-xa$, direct matrix multiplication gives
\begin{equation}\label{eq:polynomial-Gauss}
 n^+(x)n^-\!\left(\frac a{1-xa}\right)\mathbf p^j
 =n^-(a)\mathbf p^j
        \operatorname{diag}(c^{-1},c)n^+(2^jxc).
\end{equation}
Define the $V$-valued polynomial
\[
 P_j(x)=\alpha^{-j}\sum_{a\bmod2^j}
              \prod_v(1-x_va_v)^{\delta_v}\,
                    n^-(a)\mathbf p^jf.
\]
Its multidegree is bounded by~$\delta$.  Use a $G_2$-invariant
nonarchimedean norm on~$V$.  The torus eigenrelation,
\eqref{eq:polynomial-Gauss}, and the ultrametric inequality give
\begin{equation}\label{eq:polynomial-uniform-error}
 \sup_{x\in X}\|n^+(x)f-P_j(x)\|
 \le \sup_{y\in(2^{j+m}\ZZ_2)^{[F:\QQ]}}
                         \|n^+(y)f-f\|\longrightarrow0.
\end{equation}
There is no factor for the number of residue classes: the norm of a
sum is bounded by the largest norm of a summand.  The unit eigenvalue
and unitarity also introduce no growing factor.  Continuity at the
identity proves the final convergence.

For each~$v$, choose $\delta_v+1$ distinct points in $2^m\ZZ_2$.
Evaluation on their product grid and inversion of the resulting
Vandermonde matrices recover all coefficients of a polynomial of
multidegree at most~$\delta$, continuously in its values.  This
polynomial space is therefore closed in $C(X,V)$.
Equation~\eqref{eq:polynomial-uniform-error} implies
\begin{equation}\label{eq:upper-orbit-polynomial}
                n^+(x)f=\sum_{0\le\nu\le\delta}x^\nu v_\nu.
\end{equation}

For completeness, we identify an algebraic representation containing
the resulting orbit.  Put
$H=\prod_v(1+2^mM_2(\ZZ_2))$ and write
$h_v=\left(\begin{smallmatrix}A_v&B_v\\C_v&D_v\end{smallmatrix}\right)$.
Its Gauss decomposition is
\[
 h_v=n^+(B_v/D_v)
       \operatorname{diag}(\det(h_v)/D_v,D_v)n^-(C_v/D_v).
\]
Every factor lies in the neighborhoods already considered.  Hence
\begin{equation}\label{eq:algebraic-orbit-coefficients}
 hf=\sum_{0\le\nu\le\delta}
       \prod_v\det(h_v)^{e_{1,v}}
               B_v^{\nu_v}D_v^{\delta_v-\nu_v}\,v_\nu.
\end{equation}
These are regular algebraic functions on~$G_2$, also when some
$e_{1,v}$ is negative, since powers of the determinant are invertible
in its coordinate ring.

Let $z_0=\bigotimes_ve_2^{\delta_v}\in A$, with the determinant
factors understood.  In the binomially scaled monomial basis $E_\nu$
of~$A$, one has $n^+(x)z_0=\sum_\nu x^\nu E_\nu$.
Define $\Theta:A\to V$ by $\Theta(E_\nu)=v_\nu$.
Formula~\eqref{eq:algebraic-orbit-coefficients} says
$\Theta(hz_0)=hf$ for $h\in H$.
The vectors $kz_0$, $k\in H$, span~$A$, by the same upper-unipotent
interpolation.  Thus
\[
 \Theta(hkz_0)=hkf=h\Theta(kz_0),\qquad h,k\in H,
\]
proves that $\Theta$ is $H$-equivariant.  Since $H$ is Zariski dense,
its invariant subspace $\ker\Theta\subset A$ is an algebraic
subrepresentation.  The representation $A$ is irreducible and
$\Theta(z_0)=f\ne0$, so $\Theta$ is injective.  Its image contains~$f$
and has algebraic type~$A$.  This proves local
algebraicity without any assertion that algebraic vectors commute
with specialization.
\end{proof}

The vector of Lemma~\ref{lem:eigenvector-specialization} meets these
hypotheses by \eqref{eq:nonnegative-algebraic-degree}.  It is therefore
a nonzero locally algebraic vector with the target eigensystem.

\subsection{Completion of the proof over \texorpdfstring{$\QQ$}{Q}}

If $r|_{\GQtwo}$ is absolutely irreducible, the preceding ordinary
argument is unnecessary.  Proposition~\ref{prop:block-comparison}
and the local correspondence give nonzero locally algebraic vectors
in the target eigenmodule: for an absolutely irreducible parameter,
regular de Rham is exactly the required condition
\cite[Theorems~1.1 and~1.3]{CDP2014},
\cite[Theorem~0.20]{Colmez2010}.  The all-prime correspondence covers
the scalar residual case at~$2$.  Since $2$ splits completely in~$F$,
this applies at every dyadic factor.
The parameter in this application is $r\otimes\eps$.  We translate
the opposite Hodge--Tate sign convention in these local correspondence
sources using cyclotomic exponents, obtaining the algebraic factor fixed
in Proposition~\ref{prop:block-comparison}.

In either local case, the definite-forms comparison of
Proposition~\ref{prop:completed-forms} now gives a regular algebraic
automorphic eigensystem on the definite quaternion algebra over~$F$.
It cannot be a reduced-norm character, since its global Galois
representation is absolutely irreducible.  Jacquet--Langlands therefore
gives a regular algebraic Hilbert cusp form
\cite[Theorem~14.4]{JacquetLanglands1970}.  Its Galois representation
exists at~$2$ by \cite[Theorem~2.1.1]{BLGGT2014} and agrees with
$r|_{G_F}$ by the unramified Hecke identities and Chebotarev.

Finally choose a tower of prime-degree cyclic extensions from~$\QQ$ to a solvable
Galois closure of~$F$, replacing $F$ by that closure if needed.
At each stage the Hilbert representation upstairs is invariant under
the cyclic Galois group: its conjugates have the same unramified
eigensystems because the Galois representation extends downstairs, so
strong multiplicity one identifies them
\cite[Chapter~3, Lemma~3.1]{Langlands1980}.  Cyclic descent supplies a
cuspidal representation downstairs, unique up to characters of that
cyclic extension
\cite[Chapter~III, Theorems~4.2(d) and~5.1]{ArthurClozel1989}; see also
the original degree-two construction in
\cite[Chapter~11, Lemma~11.3 and Proposition~11.4]{Langlands1980}.
When needed, apply the cited unitary formulation after twisting by a
power of the idelic norm, and undo that twist after descent.  Every
real place splits in the totally real extension, so the corresponding
local extension is $\mathbb R/\mathbb R$; archimedean base change
therefore preserves the regular algebraic real-place type.
The Galois representation of the descended packet and the given
representation become isomorphic on
restriction upstairs; absolute irreducibility of that restriction
implies that they differ by a character of the cyclic quotient.
Twist the descended automorphic representation by the inverse finite
Hecke character to correct this difference, and continue down the
tower.  Over~$\QQ$ the resulting representation is attached to a
classical cusp eigenform, with the Tate twist dictated by the original
weights.  This proves Theorem~\ref{thm:main}.
